\noindent This chapter covers all 37 supplied catalog records, from
6~September~1982 to 18~April~1983: 36 mathematical records and the year heading
579. The occurrence index contains 335 items. Catalog page ranges are
retained as supplied; item locators follow the actual handwriting. On the two
April meeting dates, the source puts 594 before 593 and 596 before 595.
\GSpage{P001115}, although included in the input packet, begins Deligne's
5~September~1983 talk and belongs to the following academic year. The present
mathematical coverage therefore ends at \GSpage{P001114}.

\paragraph{The year's connected lines.}
Affine Lie algebras, their Heisenberg and vertex realizations, the Virasoro
algebra, and semi-infinite constructions connect Kirillov, Fuchs, Feigin,
Neretin, and the string lectures. Zamolodchikov and Polyakov move between
classical constraints, random surfaces, conformal symmetry, and Liouville
theory; Schwarz and Kobzarev supply the gauge-theoretic and defect background.
The companion keeps formal path integrals distinct from later mathematical
constructions. In particular, algebraic quantization of cyclic string charges
has an established existence result
\citep{enriquezhalbout2010quantization,bahnsmeinecke2011string}; it is not
presented as a newly unresolved existence question.

The geometric line runs through orbital integrals, canonical relations,
moment maps, and localization. Beilinson's disk and gluing discussion has an
unusually strong documentary match: his later paper explicitly follows
S.~I.~Gelfand's 1982 notes \citep[introduction]{beilinson1987glue}.
Manin's superflags and Gelfand's continuous Grassmannians have more uneven
levels of formulation. In the spring, the finite-gap and Hamiltonian line
joins Krichever's perturbation theory, Novikov's many-valued actions,
Semenov-Tian-Shansky's factorization, and Wilson's Grassmannian construction
of integrable hierarchies \citep{krichever1983hessians,semenov1983classical,segalwilson1985loop}.
Wodzicki's spectral zeta functions and boundary index theory create a bridge
to the final arithmetic discussions without identifying spectral and
arithmetic zeta functions indiscriminately.

\paragraph{What the classifications mean.}
The counts describe occurrences, including recurring exposition, rather than
independent discoveries. A publication class and literal correctness are
separate: an old theorem can be copied with a wrong sign, and a contemporary
paper can repair a false notebook formula.
\begin{center}
\begin{tabular}{@{}lrrrrrrr@{}}
\toprule
Class & A & B & C & D & E & F & Total\\
\midrule
Items & 47 & 29 & 3 & 0 & 214 & 42 & 335\\
\bottomrule
\end{tabular}
\end{center}
Here A denotes near-contemporary publication by the speaker; B, a later
established development by others; C, a specified partial resolution; D,
documented lack of development; E, prior knowledge; and F, insufficient
identification. No class D is assigned merely because a search failed.
The separate assessment axes and source audit distinguish direct reading
of a result from bibliographic metadata checked.

\paragraph{Questions and bounded development.}
The problem index has 19 entries, including exact recurrences.
The genuine remaining tasks include the global four-unitary free energy
(GP-1982-02), general continuation for curve zeta functions (GP-1982-42), and
actual automorphy of arbitrary curve cohomology (GP-1982-43).
For the matrix integral, the development proves finite-rank formulas and
uses the small-potential theorem of
\citet[Theorem~0.1]{collinsguionnetmaurel2009asymptotics} to justify
\(F_4(\beta)=-4\beta+6\beta^2+O(\beta^3)\) locally, in the notebook's
negative-exponent convention. An independent exact index contraction checks
the coefficients. This does not evaluate the integral at every coupling.
In arithmetic, the genus-two continuation theorem over totally real fields
\citep[Theorems~1.1.2--1.1.3]{boxercalegarigeepilloni2021surfaces}
must be distinguished from actual modularity with additional hypotheses.
The broad remarks on continuous Grassmannian topology, arbitrary abelian
theta functions, and group-free localization remain explicitly underspecified.

\paragraph{Continuities across years.}
The earlier occurrence index supplies exact returns to compact-group weight
multiplicities (GP-1981-41), eta regularity (GP-1981-43), and spectral
constructions without self-adjointness (GP-1981-44). The finite-gap and KP
lectures continue the 1973--77 sequence; the superspace discussion continues
1978-455 and 1981-537. These subject links are kept distinct from an exact
recurrence of a question. All present-status judgments are dated
November~2026; uncertain readings and unresolved identifications remain
visible rather than being completed by historical guesswork.
